\chapter{Evolution equations and comparison}\label{sec:mcf-evolution}
We work with smooth hypersurfaces in Euclidean $\mathbb R^{n+1}$, $n\geq2$, using $H=\operatorname{div}_\Sigma\nu$ and $\vec H=-H\nu$. Unless specified otherwise, a flow is parametrized by $\partial_tF=-H\nu$. The metric and covariant derivatives below are those of the evolving hypersurface. Ambient curvature would add terms to these formulas; the present statements are Euclidean. Sources are Huisken--Polden \cite{HP99}, Sections 2--4, and Huisken \cite{Hui90}, Section 1.

\begin{definition}[Scalar second fundamental form]\label{def:mcf-scalar-second-form}
For a chosen unit normal $\nu$ and local coordinate vectors $e_i=\partial_iF$, put
\[
 h_{ij}=\langle D_{e_i}\nu,e_j\rangle,
 \qquad A_i{}^j=g^{jk}h_{ik},\qquad H=\operatorname{tr}A.
\]
Thus $h$ is positive on the outward-oriented round sphere, and the vector second fundamental form of Definition \ref{def:smooth-minimal} is $B_{ij}=-h_{ij}\nu$. We write $|A|^2=h_{ij}h^{ij}=|B|^2$. Raising an index in a time derivative requires differentiating $g^{ij}$ as well. This convention agrees with \cite{HP99}, Section 2.
\end{definition}

\begin{proposition}[Euclidean normal-variation formulas]\label{prop:normal-variation-formulas}
For a smooth variation with velocity $\partial_tF=f\nu$, one has
\[
 \partial_tg_{ij}=2fh_{ij},\qquad
 \partial_tg^{ij}=-2fh^{ij},\qquad
 \partial_t d\mu=fH\,d\mu,\qquad \partial_t\nu=-\nabla f,
\]
\[
 \partial_th_{ij}=-\nabla_i\nabla_jf+fh_i{}^kh_{kj},
 \qquad \partial_tH=-\Delta f-f|A|^2.
\]
These identities do not require the initial hypersurface to be minimal. Compare \cite{HP99}, Theorem 3.2, p.~52, specialized to Euclidean ambient space.
\end{proposition}
\begin{proof}
Differentiate $g_{ij}=\langle\partial_iF,\partial_jF\rangle$ and use $\partial_i(f\nu)=f_i\nu+fD_i\nu$. Differentiate $g^{ik}g_{kj}=\delta^i_j$ to obtain the inverse-metric formula; the determinant derivative gives the measure formula. Differentiating $\langle\nu,\nu\rangle=1$ and $\langle\nu,e_i\rangle=0$ shows that $\partial_t\nu$ is tangent with components $-f_i$.

At a point choose coordinates with $g_{ij}=\delta_{ij}$ and vanishing Christoffel symbols. Differentiate $h_{ij}=\langle D_i\nu,e_j\rangle$. The derivative of $\nu$ contributes $-\nabla_i\nabla_j f$, while the derivative of $e_j$ contributes $f\langle D_i\nu,D_j\nu\rangle$. This proves the formula for $h$. Finally differentiate $H=g^{ij}h_{ij}$: the inverse-metric contribution is $-2f|A|^2$, leaving $-\Delta f-f|A|^2$ after tracing the $h$ formula.
\end{proof}

\begin{theorem}[Euclidean Simons identity]\label{thm:euclidean-simons}
For a smooth Euclidean hypersurface with the convention in Definition \ref{def:mcf-scalar-second-form},
\[
 \Delta h_{ij}=\nabla_i\nabla_jH+Hh_i{}^kh_{kj}-|A|^2h_{ij}.
\]
This is the Euclidean specialization of \cite{HP99}, Corollary 2.2, p.~49. The formula is valid without a minimality or convexity hypothesis.
\end{theorem}
\begin{proof}
Apply Theorem \ref{thm:spaceform-simons} with $c=0$. Its proof gives the Gauss--Codazzi commutation explicitly, so no external curvature-commutation input is needed for the evolution identities below.
\end{proof}

\begin{proposition}[Curvature evolution under mean curvature flow]\label{prop:mcf-curvature-evolution}
Along a smooth Euclidean hypersurface flow,
\[
 (\partial_t-\Delta)A_i{}^j=|A|^2A_i{}^j,
 \qquad (\partial_t-\Delta)H=|A|^2H,
\]
\[
 (\partial_t-\Delta)|A|^2=-2|\nabla A|^2+2|A|^4.
\]
The mixed-index equation differs from the equation for $h_{ij}$, because the metric evolves. See \cite{HP99}, Section 3, and \cite{Hui90}, p.~286.
\end{proposition}
\begin{proof}
Put $f=-H$ in Proposition \ref{prop:normal-variation-formulas}. Theorem \ref{thm:euclidean-simons} gives
\[
 \partial_th_{ij}=\Delta h_{ij}+|A|^2h_{ij}-2Hh_i{}^kh_{kj}.
\]
Since $\partial_tg^{ij}=2Hh^{ij}$, raising an index cancels the last term. Trace the resulting mixed-index equation to obtain the equation for $H$. Finally $|A|^2=\operatorname{tr}(A^2)$, so the product rule and metric compatibility give
$\Delta\operatorname{tr}(A^2)=2\operatorname{tr}(A\Delta A)+2|\nabla A|^2$.
Taking its time derivative proves the last identity.
\end{proof}

\begin{corollary}[Preservation of nonnegative mean curvature]\label{cor:mcf-mean-convexity}
Let a smooth flow have compact connected parameter manifold without boundary and a continuous choice of normal. If $H\geq0$ initially, then $H\geq0$ at every later smooth time. If initially $H$ is not identically zero, then $H>0$ for every positive smooth time.
\end{corollary}
\begin{proof}
On each compact time interval contained in the smooth existence interval, $|A|^2$ is bounded. Apply the weak and strong parabolic maximum principles to $\partial_tH=\Delta H+|A|^2H$, multiplying by $e^{-Ct}$ when needed to make the zeroth-order coefficient nonpositive. Connectedness is used for the strict conclusion. This is an application of the scalar maximum principle, as in \cite{Hui90}, Section 4; it does not assert preservation of every curvature pinching condition.
\end{proof}

\begin{proposition}[Position and squared-distance equations]\label{prop:mcf-position-heat}
For a normal parametrization of a smooth Euclidean mean curvature flow and a fixed $a\in\mathbb R^{n+1}$,
\[
 \partial_tF=\Delta_{\Sigma_t}F,\qquad
 (\partial_t-\Delta_{\Sigma_t})|F-a|^2=-2n.
\]
Every fixed linear coordinate of $F$ solves the scalar heat equation for the evolving metric.
\end{proposition}
\begin{proof}
The normal component of $D_ie_j$ is $-h_{ij}\nu$. Tracing after removing the tangential Christoffel term gives $\Delta F=-H\nu=\partial_tF$. The product rule gives
$\Delta|F-a|^2=2\langle F-a,\Delta F\rangle+2|\nabla F|^2$.
An orthonormal tangent frame yields $|\nabla F|^2=n$, proving the second identity. Compare the Euclidean identities underlying \cite{Hui90}, Sections 1--2.
\end{proof}

\begin{corollary}[Confinement and finite lifetime]\label{cor:mcf-confinement}
If a nonempty smooth closed flow starts inside $\overline{B_{R_0}(a)}$, then at every smooth time
\[
 \max_{\Sigma_t}|x-a|^2\leq R_0^2-2nt.
\]
Consequently its smooth maximal existence time is at most $R_0^2/(2n)$. Each linear coordinate also stays between its initial minimum and maximum, so the flow stays in the convex hull of its initial image.
\end{corollary}
\begin{proof}
Apply the maximum principle to $|F-a|^2+2nt$ and, separately, to each linear coordinate in Proposition \ref{prop:mcf-position-heat}. A nonempty immersed positive-dimensional hypersurface cannot have negative squared distance or consist of the single point $a$. Intersecting the preserved half-spaces gives the convex-hull assertion.
\end{proof}

\begin{proposition}[Avoidance for two closed hypersurface flows]\label{prop:mcf-avoidance}
Let $F_i:M_i^n\times[0,T)\to\mathbb R^{n+1}$, $i=1,2$, be smooth mean curvature flows of compact manifolds without boundary. If their initial images are disjoint, then their distance is nondecreasing on $[0,T)$. In particular their images remain disjoint while both flows are smooth.
\end{proposition}
\begin{proof}
Use normal parametrizations. Until a possible first contact, choose a pair realizing the positive minimum of $|F_1(p,t)-F_2(q,t)|^2$, and write $d=F_1-F_2$. The first derivative test makes $d$ normal to both tangent spaces. Since both are hypersurfaces, the tangent spaces coincide with $d^\perp$. Choose matching orthonormal tangent vectors $e_1,\ldots,e_n$ at the pair, and vary the two points simultaneously along geodesics with these initial velocities. The first derivatives of $d$ cancel. Summing the second derivative test yields
\[
 0\leq2\langle d,\vec H_1-\vec H_2\rangle
 =\partial_t|F_1-F_2|^2.
\]
On compact time intervals the minimum is locally Lipschitz; its derivative at almost every time is bounded below by the least time derivative among its minimizing pairs. Hence the squared distance cannot decrease. Its initial positive value rules out first contact. This is the closest-pair argument illustrated for curves in \cite{WhiMCF15}, Theorem 2.2; the common-tangent-space step explains the hypersurface hypothesis here.
\end{proof}

\begin{theorem}[Short-time existence and the curvature continuation criterion]\label{thm:mcf-existence-continuation}
A smooth closed immersed hypersurface in $\mathbb R^{n+1}$ admits a smooth mean curvature flow on a nonempty time interval. There is a finite maximal smooth time $T>0$, and
\[
 \limsup_{t\uparrow T}\max_{\Sigma_t}|A|=\infty.
\]
Equivalently, a uniform bound for $|A|$ up to a finite endpoint permits smooth extension beyond it. Short-time existence is the mean-curvature specialization of \cite{HP99}, Theorem 3.1 and its subsequent example, pp.~51--52. The maximal-time and continuation statement is \cite{Hui90}, Theorem 1.1, p.~286. The parabolic existence and derivative-estimate proofs remain external. This theorem does not furnish a continuation through a curvature singularity in the smooth category.
\end{theorem}

\begin{proposition}[Curvature ODE bound and a lower singularity rate]\label{prop:mcf-curvature-rate}
For a smooth closed flow, put $U(t)=\max_{\Sigma_t}|A|^2$. If $t\geq s$ and $2U(s)(t-s)<1$, then
\[
 U(t)\leq\frac{U(s)}{1-2U(s)(t-s)}.
\]
If $T$ is the finite maximal smooth time, then
\[
 U(s)\geq\frac{1}{2(T-s)}\quad(0\leq s<T).
\]
Thus the second inequality is a lower blow-up rate, not a Type I upper bound. Compare \cite{Hui90}, Lemma 1.2.
\end{proposition}
\begin{proof}
At a maximum, Proposition \ref{prop:mcf-curvature-evolution} gives the upper differential inequality $U'\leq2U^2$, valid almost everywhere on compact smooth time intervals. Compare with the solution $y'=2y^2$, $y(s)=U(s)$, to obtain the first estimate. If $2U(s)(T-s)<1$, the same estimate bounds curvature up to $T$, contradicting Theorem \ref{thm:mcf-existence-continuation}. This proves the second assertion without presuming that the blow-up has a particular model.
\end{proof}

\begin{definition}[Type I and Type II finite-time singularities]\label{def:mcf-singularity-types}
For a closed flow with finite maximal smooth time $T$, the singularity is Type I if
\[
 \sup_{0\leq t<T}(T-t)\max_{\Sigma_t}|A|^2<\infty.
\]
Otherwise it is Type II. This is a global classification of the rate at the first singular time, and does not by itself specify the model at every singular point. Compare the upper-rate assumption (2) in \cite{Hui90}, p.~286, and the rescaling discussion in \cite{HP99}, Section 4.
\end{definition}

\begin{proposition}[Parabolic rescaling]\label{prop:mcf-parabolic-rescaling}
Fix $\lambda>0$ and $(x_0,t_0)$. Whenever the original time is defined, set
\[
 \widetilde F(p,s)=\lambda\bigl(F(p,t_0+\lambda^{-2}s)-x_0\bigr).
\]
Then $\widetilde F$ is again mean curvature flow. Its area and curvature satisfy
\[
 d\widetilde\mu=\lambda^n d\mu,\qquad
 |\widetilde A|=\lambda^{-1}|A|,\qquad
 \widetilde H=\lambda^{-1}H.
\]
In particular $(T-t)|A|^2$ is unchanged by the corresponding change of space and time scales. See \cite{WhiMCF15}, Section 6.
\end{proposition}
\begin{proof}
Spatial dilation changes the induced metric to $\lambda^2g$ and the covariant second fundamental form to $\lambda h$. Consequently the shape operator scales by $\lambda^{-1}$, and the area element by $\lambda^n$. Differentiating the displayed parametrization in $s$ gives $\lambda^{-1}\partial_tF=\lambda^{-1}\vec H$, which is the new mean curvature vector. The time-to-endpoint scales by $\lambda^2$, proving the last claim.
\end{proof}

\begin{proposition}[The graphical mean curvature equation]\label{prop:graphical-mcf}
The family of graphs $F(x,t)=(x,u(x,t))$ in $\mathbb R^{n+1}$ moves by mean curvature in its normal velocity precisely when
\[
 \partial_tu=\sqrt{1+|Du|^2}\,\operatorname{div}\left(\frac{Du}{\sqrt{1+|Du|^2}}\right)
 =\left(\delta_{ij}-\frac{u_i u_j}{1+|Du|^2}\right)u_{ij}.
\]
On a region with $|Du|\leq L$, the spatial coefficient eigenvalues lie between $(1+L^2)^{-1}$ and $1$. The parametrization need not have purely normal velocity. See \cite{WhiMCF15}, Section 3.1, p.~8.
\end{proposition}
\begin{proof}
The upward normal is $\nu=(-Du,1)/\sqrt{1+|Du|^2}$, so the normal speed is $u_t/\sqrt{1+|Du|^2}$. With our scalar convention, $H=-\operatorname{div}(Du/\sqrt{1+|Du|^2})$. Equating normal speed with $-H$ proves the divergence expression; expanding it gives the second expression. Its eigenvalues were computed in Proposition \ref{prop:graph-ellipticity}.
\end{proof}

\begin{proposition}[Comparison for graphical flows]\label{prop:graphical-mcf-comparison}
Let $\Omega\subset\mathbb R^n$ be bounded with smooth boundary. Suppose $u_1,u_2$ are smooth solutions of Proposition \ref{prop:graphical-mcf} on $\overline\Omega\times[0,T]$. If $u_1\leq u_2$ on $\overline\Omega\times\{0\}$ and on $\partial\Omega\times[0,T]$, then $u_1\leq u_2$ throughout the cylinder.
\end{proposition}
\begin{proof}
Write $w=u_2-u_1$ and $a^{ij}(p)=\delta_{ij}-p_ip_j/(1+|p|^2)$. Subtracting the equations gives
\[
 w_t=a^{ij}(Du_2)w_{ij}+b^kw_k,
 \qquad b^k=\int_0^1\partial_{p_k}a^{ij}(Du_1+\theta Dw)\,d\theta\,(u_1)_{ij}.
\]
The smoothness and compactness assumptions bound the gradients and $b$, giving uniform parabolicity. The parabolic maximum principle implies $w\geq0$. This states explicitly the boundary and regularity conditions for the comparison argument in \cite{WhiMCF15}, Section 3.1, pp.~8--9.
\end{proof}
